\documentclass[11pt]{article}
\usepackage[margin=1in]{geometry}
\geometry{letterpaper}

\usepackage{amsmath}
\usepackage{amssymb}
\usepackage{tikz-cd}

%%% Theorems and references %%%
\usepackage[amsmath,thmmarks]{ntheorem}
\usepackage{hyperref}
\usepackage{cleveref}

\theoremstyle{change}

\newtheorem{definition}[equation]{Definition}
\newtheorem{theorem}[equation]{Theorem}
\newtheorem{prop}[equation]{Proposition}
\newtheorem{lemma}[equation]{Lemma}
\newtheorem{cor}[equation]{Corollary}
\newtheorem{conj}[equation]{Conjecture}
\newtheorem{problem}[equation]{Problem}
\newtheorem{question}[equation]{Question}

\theorembodyfont{\upshape}
\theoremsymbol{\ensuremath{\Diamond}}
\newtheorem{remark}[equation]{Remark}
\newtheorem{convention}[equation]{Convention}

\theoremstyle{nonumberplain}

\theoremsymbol{\ensuremath{\Box}}
\newtheorem{proof}{Proof}

\qedsymbol{\ensuremath{\Box}}

\creflabelformat{equation}{#2(#1)#3}

\crefname{equation}{equation}{equations}
\crefname{definition}{definition}{definitions}
\crefname{theorem}{Theorem}{Theorems}
\crefname{prop}{proposition}{propositions}
\crefname{lemma}{lemma}{lemmas}
\crefname{cor}{corollary}{corollaries}
\crefname{conj}{conjecture}{conjectures}
\crefname{problem}{problem}{problems}
\crefname{question}{question}{questions}
\crefname{remark}{remark}{remarks}
\crefname{convention}{convention}{conventions}
\crefname{section}{Section}{Sections}
\crefname{subsection}{Section}{Sections}

\crefformat{equation}{#2equation~(#1)#3}
\crefformat{definition}{#2definition~#1#3}
\crefformat{theorem}{#2Theorem~#1#3}
\crefformat{prop}{#2proposition~#1#3}
\crefformat{lemma}{#2lemma~#1#3}
\crefformat{cor}{#2corollary~#1#3}
\crefformat{conj}{#2conjecture~#1#3}
\crefformat{problem}{#2problem~#1#3}
\crefformat{question}{#2question~#1#3}
\crefformat{remark}{#2remark~#1#3}
\crefformat{convention}{#2convention~#1#3}
\crefformat{section}{#2Section~#1#3}
\crefformat{subsection}{#2Section~#1#3}

\Crefformat{equation}{#2Equation~(#1)#3}
\Crefformat{definition}{#2Definition~#1#3}
\Crefformat{theorem}{#2Theorem~#1#3}
\Crefformat{prop}{#2Proposition~#1#3}
\Crefformat{lemma}{#2Lemma~#1#3}
\Crefformat{cor}{#2Corollary~#1#3}
\Crefformat{conj}{#2Conjecture~#1#3}
\Crefformat{problem}{#2Problem~#1#3}
\Crefformat{question}{#2Question~#1#3}
\Crefformat{remark}{#2Remark~#1#3}
\Crefformat{convention}{#2Convention~#1#3}
\Crefformat{section}{#2Section~#1#3}
\Crefformat{subsection}{#2Section~#1#3}

%%% Letters, Symbols, Words %%%

\newcommand\NN{{\mathbb N}}
\newcommand\QQ{{\mathbb Q}}
\renewcommand\AA{{\mathbb A}}
\newcommand\GG{{\mathbb G}}

\newcommand\cJ{{\mathcal J}}
\newcommand\cW{{\mathcal W}}
\newcommand\cR{{\mathcal R}}
\newcommand\cX{{\mathcal X}}
\newcommand\cY{{\mathcal Y}}
\newcommand\cD{{\mathcal D}}
\newcommand\cK{{\mathcal K}}
\newcommand\fE{{\mathfrak E}}

\newcommand\mono{\hookrightarrow}
\newcommand\sminus{\smallsetminus}
\newcommand\st{{\textrm{ s.t.\ }}}
\newcommand\setof[1]{\{ #1 \}}
\newcommand\rev{{\mathrm{rev}}}
\newcommand\FI{{\mathrm{FI}}}

\parskip 1ex

\DeclareMathOperator{\End}{End}
\DeclareMathOperator{\Aut}{Aut}
\DeclareMathOperator{\Cent}{Cent}
\DeclareMathOperator{\Frac}{Frac}
\DeclareMathOperator{\GL}{GL}
\DeclareMathOperator{\Iso}{Iso}
\DeclareMathOperator{\id}{id}
\DeclareMathOperator{\im}{im}
\DeclareMathOperator{\Gpd}{Gpd}
\DeclareMathOperator{\Span}{Span}
\DeclareMathOperator{\Mon}{Mon}
\DeclareMathOperator{\Ind}{Ind}
\DeclareMathOperator{\Constr}{Constr}
\DeclareMathOperator{\Fun}{Fun}
\DeclareMathOperator{\stab}{stab}

\title{The Jacobian--Dixmier Monoid}

\begin{document}
\maketitle

\begin{abstract}
The explicit counterexamples to the Jacobian conjecture in dimension $3$ \cite{Alpoge,Speyer,Tao,fam} settle $JC_n$ and $DC_n$ for all $n \geq 3$, but for the present purposes they are only seeds. The object of study is the structure they seed: for each $n$ the action groupoid of $\Aut(W_n) \times \Aut(W_n)^{op}$ on $\End(W_n)$, with its composition law, its isotropy, and the coordinate stabilizations $n \to n+1$. We organize this as a monoid in spans of groupoids, graded by a quantum degree $\deg \psi = \dim_{E_\psi} \cD_n$ that we prove agrees with the generic fibre degree on Bass--Connell--Wright lifts, and we package the stabilizations as a functor out of the category $\FI$ of finite sets and injections. The whole situation is then a single diagram $\cR^J \Rightarrow \cJ^\rev \Rightarrow \cW$ of $\FI$-indexed graded monoids in spans, where $\cR^J$ is the part generated by the known Keller maps. Two questions remain: whether $\cR^J \Rightarrow \cJ^\rev$ is \emph{enough} (every Keller map is a word in the seeds and automorphisms) and whether $\cR^J \Rightarrow \cW$ is enough (every endomorphism of $W_n$ is a word in lifted seeds and automorphisms of $W_n$). With the seeds restricted to Alp\"oge's map alone, the second question has a negative answer as soon as a Keller map of generic fibre degree $4$ is certified; with the full Gallagher family as seeds, degree separates nothing and finer invariants are needed.
\end{abstract}

\begin{convention}\label{conv}
Throughout, $k$ is an algebraically closed field of characteristic $0$. Characteristic $0$ is used for simplicity of the Weyl algebra, separability, and Keller's theorem; algebraic closure is used so that images of constructible sets of $k$-points are constructible (Chevalley). The motivating case is $k = \mathbb{C}$, as the coefficient field of the quantized polynomial observables on a $2n$-dimensional phase space; nothing is specific to it. No further structure on $k$ (real structure, Galois action) is used.
\end{convention}

\section{Jacobian and Dixmier background}

\begin{definition}[$JC_n$]
Any polynomial endomorphism $F$ of $\AA^n_k$ with Jacobian determinant $1$ is a polynomial automorphism. We call a polynomial endomorphism with constant nonzero Jacobian determinant a \emph{Keller map}.
\end{definition}

\begin{theorem}[Normalization; Connell -- van den Dries\label{vC}]
Let $F$ be a Keller map with $\det JF = \lambda \in k^\times$.
\begin{enumerate}
\item For $L \in \GL_n(k)$, $L \circ F$ is a Keller map with determinant $\det(L)\lambda$, and is an automorphism iff $F$ is. Taking $L = \mathrm{diag}(\lambda^{-1}, 1, \dots, 1)$, every Keller map normalizes to determinant $1$.
\item \cite{CvdD} The truth of $JC_n$ is independent of the field of characteristic $0$.
\end{enumerate}
\end{theorem}

\begin{definition}[$d(F)$]
For a Keller map $F$ of $\AA^n$, the field extension $k(x_1, \dots, x_n) / k(F_1, \dots, F_n)$ is finite because $JF$ is invertible. Set $d(F) := [k(x) : k(F)]$, the \emph{degree} of $F$. It is the number of geometric points in a general fibre. It is not the polynomial degree.
\end{definition}

\begin{theorem}[Keller\label{keller}]
A Keller map with $d(F) = 1$ is a polynomial automorphism.
\end{theorem}

\begin{theorem}[Campbell \cite{Campbell}\label{campbell}]
If $k(x)/k(F)$ is a Galois extension then $F$ is a polynomial automorphism. In particular, for a Keller map that is not an automorphism, $k(x)/k(F)$ is not Galois.
\end{theorem}

\begin{definition}[$W_n$]
The $n$-th Weyl algebra over $k$ is the unital associative $k$-algebra with $2n$ generators $q_i, \partial_i$ and relations $[q_i, \partial_i] = -1$, all other pairs commuting. It is the algebra of polynomial differential operators on $\AA^n$, and the quantization of the polynomial functions on the $2n$-dimensional phase space $T^*\AA^n$. It carries no further structure.
\end{definition}

$W_n$ is simple, so every endomorphism of $W_n$ is injective, and $W_n$ is a Noetherian domain, so it has a skew field of fractions $\cD_n$.

\begin{definition}[$DC_n$]
Every endomorphism of $W_n$ is an automorphism, i.e.\ surjective.
\end{definition}

\begin{theorem}[Belov-Kanel, Kontsevich \cite{BKK}]
$DC_n \implies JC_n$ and $JC_{2n} \implies DC_n$.
\end{theorem}

Only the first implication is used here, in the constructive form of \cref{lift}. The second is proved by reduction modulo $p$ with choices made for each endomorphism separately (spreading out, an ultrafilter); it does not respect composition or stabilization and plays no structural role below.

\section{The seeds\label{seed}}

\subsection{Alp\"oge's map}

The counterexample announced by Alp\"oge on 19 July 2026 \cite{Alpoge,Speyer} is $F = (a,b,c): \AA^3 \to \AA^3$ with
\begin{align*}
a&=(1+xy)^3z+y^2(1+xy)(4+3xy),\\
b&=y+3x(1+xy)^2z+3xy^2(4+3xy),\\
c&=2x-3x^2y-x^3z.
\end{align*}
Its Jacobian determinant is $-2$ and $d(F) = 3$. The fibre over $(0,3,0)$ is
\begin{equation*}
\left( -\tfrac{2}{3},\ 3,\ 18 \right), \qquad \left( 0,\ 3,\ -36 \right), \qquad \left( \tfrac{2}{3},\ -\tfrac{3}{2},\ \tfrac{45}{4} \right).
\end{equation*}
By \cref{vC}, $F_1 := (a, b, -c/2)$ has determinant $1$ and the same fibres. The components are homogeneous of weights $2, 1, -1$ for the grading $\deg(x,y,z) = (-1,1,2)$ \cite{Speyer}.

\subsection{Gallagher's family}

Gallagher \cite{fam} constructs Keller maps of $\AA^3$ of every degree $d \geq 3$ (we use \emph{degree} for $d(F)$ throughout). Given a polynomial $p$ with $p(0) = 0$, a constant $c \neq 0$ with $p(1) = -c$ and $\int_0^1 p = 0$, define $q$ by $q' = w p'/c$, $q(0) = 0$; set $\kappa = p'(1)/c$ (required $\neq -2$), $a = -(1+\kappa)/(2+\kappa)$, and let $b \neq 0$ be arbitrary. With $u = 1+xy$, $\gamma = 1 + a\,xy + b\,x^2 z$ and $w = u\gamma$,
\begin{equation*}
F_p(x,y,z) = \Big( \frac{u + q(w)/\gamma^2}{x^2},\ \frac{c + p(w)/\gamma}{x},\ x\gamma \Big)
\end{equation*}
is asserted to be a polynomial map with $\det JF_p = bc$ and $d(F_p) = \deg p + 1$. The seed $p = 2w - 3w^2$ returns Alp\"oge's map; the seeds $p_d(w) = 2w - 3w^2 + w(1-w)\big(w^{d-2} - 6/(d(d+1))\big)$ give every degree $d+1 \geq 3$. All members are weighted homogeneous for the same grading as Alp\"oge's map, since $\gamma$ and $u$ depend only on $xy$ and $x^2z$.

\begin{remark}[Status of the degree-$4$ member\label{deg4}]
For $p = w - 2w^3$, i.e.\ $u = 1+3xy$ and $\gamma = 1 - 4xy - x^2 z$, the map is
\begin{equation*}
G(x,y,z) = \Big( \frac{2u + u^2 - 3u^4\gamma^2}{x^2},\ \frac{1 + u - 2u^3\gamma^2}{x},\ x\gamma \Big).
\end{equation*}
We verified symbolically that $G$ is polynomial with component degrees $(12, 11, 4)$, that $\det JG = -6$, that $G(1,0,0) = G(-1,0,2) = (0,0,1)$, and, by Gr\"obner bases at three random rational target points, that those fibres consist of exactly $4$ points. Since $G$ is \'etale, a fibre with $4$ points gives $d(G) \geq 4$. The claim $d(G) = 4$ exactly is Gallagher's; it follows from his fibre equation $\int_0^w p = wP - cQ$, a quartic in $w$, provided $(x,y,z)$ is rational in $w$ and the target, which we have not independently certified. Everything below that depends on $d(G) = 4$ is marked as such.
\end{remark}

\begin{definition}[Seed set]
Let $S_3 \subset \End(\AA^3)$ be the set of normalized (determinant $1$) Keller maps consisting of $F_1$ and all members of Gallagher's family. For $n \geq 3$ and a coordinate injection $f : [3] \mono [n]$, $\stab_f(F)$ denotes $F$ acting on the coordinates in the image of $f$ and the identity on the others. The \emph{seeds in dimension $n$} are $S_n := \setof{ \stab_f(F) \st F \in S_3,\ f : [3] \mono [n] }$.
\end{definition}

\section{The lift\label{liftsec}}

\begin{definition}[$P_n$]
$P_n := k[x_1, \dots, x_n, y_1, \dots, y_n]$ with the Poisson bracket $\{x_i, y_j\} = \delta_{ij}$, $\{x_i,x_j\} = \{y_i,y_j\} = 0$. It is Poisson-simple: a nonzero Poisson ideal is stable under all partial derivatives, hence contains a nonzero constant.
\end{definition}

\begin{prop}[The lift\label{lift}]
Let $F$ be a Keller map of $\AA^n$ with $\det JF = 1$ and $G := (JF^{\mathsf{T}})^{-1}$, a polynomial matrix. Then:
\begin{enumerate}
\item $\Phi_F : x_i \mapsto F_i(x), \; y_j \mapsto \sum_k G_{jk}(x)\, y_k$ is a Poisson endomorphism of $P_n$;
\item $\psi_F : q_i \mapsto F_i(q), \; \partial_j \mapsto D_j := \sum_k G_{jk}(q)\, \partial_k$ is an algebra endomorphism of $W_n$;
\item $\psi_{F' \circ F} = \psi_F \circ \psi_{F'}$, and $\psi_a \in \Aut(W_n)$ for $a \in \Aut(\AA^n)$.
\end{enumerate}
\end{prop}

\begin{proof}
(1) $\{F_i, F_j\} = 0$ since the $F_i$ involve only $x$. For the mixed bracket, $\{F_i, \sum_k G_{jk} y_k\} = \sum_k G_{jk}\, \partial F_i / \partial x_k = (JF \, G^{\mathsf{T}})_{ij} = \delta_{ij}$. The bracket $\{\Phi_F(y_i), \Phi_F(y_j)\} = 0$ is the integrability identity
\begin{equation}\label{integrability}
\sum_m \left( G_{jm} \frac{\partial G_{ik}}{\partial x_m} - G_{im} \frac{\partial G_{jk}}{\partial x_m} \right) = 0,
\end{equation}
which holds because $\Phi_F$ is the cotangent lift of $F$: $F$ is everywhere a local diffeomorphism, its cotangent lift $(x,p) \mapsto (F(x), (JF(x))^{-\mathsf T} p)$ is symplectic, and being symplectic is a pointwise condition.

(2) With $[f(q), \partial_k] = -\partial f / \partial x_k$, the mixed commutators give $[\psi_F(q_i), \psi_F(\partial_j)] = -\delta_{ij}$ by the same computation. Expanding $[D_i, D_j]$ in normal form gives first-order operators whose coefficients are the left side of \cref{integrability}; no ordering corrections arise since the $G_{jk}(q)$ commute and each $D_j$ has order one.

(3) Direct from the chain rule: $J(F' \circ F)^{-\mathsf T} = (JF^{-\mathsf T})\, (JF'^{-\mathsf T} \circ F)$, and the inverse of $\psi_a$ is $\psi_{a^{-1}}$.
\end{proof}

That $\psi_F$ fails to be an automorphism when $d(F) > 1$ is deduced in \cref{reflL} from the degree of \cref{degsec}, replacing the centralizer argument of Bass--Connell--Wright.

\section{The quantum degree\label{degsec}}

\begin{definition}[$\deg \psi$\label{degdef}]
For $\psi \in \End(W_n)$ let $E_\psi \subseteq \cD_n$ be the division subring generated by $\psi(W_n)$. Set
\begin{equation*}
\deg \psi := \dim_{E_\psi} \cD_n \in \setof{1, 2, 3, \dots} \cup \setof{\infty},
\end{equation*}
the dimension of $\cD_n$ as a right $E_\psi$-vector space.
\end{definition}

Since $\psi$ is injective, it extends to a ring endomorphism of $\cD_n$ whose image is a division ring containing $\psi(W_n)$ and contained in $E_\psi$; hence $\psi(\cD_n) = E_\psi$ and $\psi : \cD_n \to E_\psi$ is an isomorphism.

\begin{prop}\label{degprops}
\begin{enumerate}
\item $\deg$ is constant on $\Aut(W_n) \times \Aut(W_n)^{op}$-orbits.
\item $\deg(\psi \varphi) = \deg \psi \cdot \deg \varphi$ (with $\infty \cdot m = \infty$).
\item $\deg \psi = 1$ for $\psi \in \Aut(W_n)$.
\item For $f : [n] \mono [m]$ a coordinate injection and $\stab_f(\psi) := \psi \otimes \id$ on $W_m = W_n \otimes W_{m-n}$, $\deg \stab_f(\psi) = \deg \psi$.
\end{enumerate}
\end{prop}

\begin{proof}
(1) An automorphism $a$ of $W_n$ extends to an automorphism of $\cD_n$, and $E_{a \psi b} = a(E_\psi)$; applying the ring isomorphism $a$ to the pair $(\cD_n, E_\psi)$ preserves dimension. (2) $E_{\psi\varphi} = \psi(E_\varphi)$, and $\psi : (\cD_n, E_\varphi) \to (E_\psi, E_{\psi\varphi})$ is an isomorphism of pairs, so by the tower law $[\cD_n : E_{\psi\varphi}] = [\cD_n : E_\psi]\,[E_\psi : E_{\psi\varphi}] = [\cD_n : E_\psi]\,[\cD_n : E_\varphi]$. (3) is clear. (4) The normal-form basis $\setof{m}$ of $W_{m-n}$ gives $W_m = \bigoplus_m W_n \cdot m$ and the same decomposition for $\psi(W_n) \otimes W_{m-n}$; the Ore localization respects the tensor decomposition and the dimension is computed on the first factor.
\end{proof}

\begin{theorem}[Degree of a lift\label{degthm}]
For a Keller map $F$ with $\det JF = 1$, $\deg \psi_F = d(F)$.
\end{theorem}

\begin{proof}
Write $K = k(F) \subset L = k(q)$, $d = [L:K]$, and let $e_1, \dots, e_d$ be a $K$-basis of $L$. Put $\psi = \psi_F$.

\emph{The $D_j$ are coordinate derivatives on $K$.} $D_j(F_i) = \sum_k G_{jk}\, \partial F_i/\partial q_k = (G \cdot JF^{\mathsf T})_{ji} = \delta_{ij}$, so $D_j$ preserves $k[F]$, acts there as $\partial/\partial F_j$, and extends to $K$.

\emph{Rational-coefficient operators.} Let $R = L[\partial_1, \dots, \partial_n] \subset \cD_n$, so $\Frac R = \cD_n$. Since $G$ is invertible over $k[q]$, the $D_j$ form an $L$-basis of the commuting derivations of $L$, so every element of $R$ is uniquely $\sum_\alpha c_\alpha D^\alpha$ with $c_\alpha \in L$.

\emph{The image generates a smaller ring.} Let $S = K[D_1, \dots, D_n] \subset R$, the elements with all $c_\alpha \in K$. By the first step $S$ is the ring of differential operators on $K$ with derivations $\partial/\partial F_j$; it is a Noetherian domain, hence Ore. $\psi(W_n) = k[F][D] \subseteq S$ and every element of $K$ is a fraction of elements of $k[F]$, so $E_\psi = \Frac S$.

\emph{$R$ is free of rank $d$ over $S$.} $R = \bigoplus_i e_i S$ as right $S$-modules: expand each $c_\alpha = \sum_i e_i c_{i\alpha}$ with $c_{i\alpha} \in K$; uniqueness follows from uniqueness of normal forms and $K$-independence of the $e_i$.

\emph{The dimension is exactly $d$.} Let $V = \sum_i e_i E_\psi \subset \cD_n$. The $e_i$ are $E_\psi$-independent: if $\sum e_i \varepsilon_i = 0$, write $\varepsilon_i = s_i t^{-1}$ over a common right denominator in $S$; then $\sum e_i s_i = 0$ in $R$, so all $s_i = 0$. $V$ is closed under left multiplication by $R$ since $r e_i \in R = \sum_j e_j S$; for nonzero $s \in S$, left multiplication by $s$ is an injective right-$E_\psi$-linear endomorphism of the finite-dimensional $V$, hence bijective, so $s^{-1} V = V$; thus $E_\psi V \subseteq V$ and $V$ is a ring. $V$ is a domain of finite dimension over $E_\psi$, so left multiplication by any nonzero $v$ is bijective and $v$ is invertible; $V$ is a division ring containing $R$, hence $V = \cD_n$. So $\dim_{E_\psi} \cD_n = d$.
\end{proof}

The same argument with coefficients on the right gives the left dimension, also $d$.

\begin{cor}[Reflection\label{reflL}]
If $\psi_F \in \Aut(W_n)$ then $F \in \Aut(\AA^n)$. In particular every Keller map that is not an automorphism, in any dimension $n \geq 3$ (Alp\"oge's map, every member of Gallagher's family, every coordinate stabilization of these), gives a non-surjective endomorphism of $W_n$, and $DC_n$ fails for all $n \geq 3$.
\end{cor}

\begin{proof}
$d(F) = \deg \psi_F = 1$ by \cref{degthm,degprops}(3), so $F$ is an automorphism by \cref{keller}.
\end{proof}

\begin{cor}
For a Keller map $F$ the two sides of \cref{degthm} are finite; the lift of a degree-$d$ Keller map has degree exactly $d$, and lifts of Keller maps of different degrees lie in different orbits.
\end{cor}

\begin{remark}[Degree $2$]
No Keller map has degree $2$: a quadratic extension is Galois, so \cref{campbell} applies. Whether $\End(W_n)$ has an element of degree $2$ is open; it would be a class invisible from the Jacobian side.
\end{remark}

\section{Monoids in spans of groupoids\label{monoid}}

\subsection{The action groupoid and its composition span}

For a monoid $M$ with unit group $U$, the action groupoid $\cX$ of $U \times U^{op}$ on $M$ has objects the elements of $M$ and morphisms $(a_1, a_2) : x \to a_1 x a_2$. Its set of isomorphism classes is the double coset set $U \backslash M / U$. Composition of $M$ does not descend to a functor $\cX \times \cX \to \cX$, because the middle automorphism must be divided out; it descends to a \emph{span}.

\begin{definition}\label{spanmon}
Let $\cX^{(2)}$ be the action groupoid of $U \times U \times U^{op}$ on $M \times M$ via $(a_1, b, a_2) \cdot (x, y) = (a_1 x b^{-1}, b y a_2)$. There are functors
\begin{equation*}
\begin{tikzcd}[column sep=large]
\cX \times \cX & \cX^{(2)} \arrow[l, "\pi"'] \arrow[r, "\mu"] & \cX
\end{tikzcd}
\end{equation*}
with $\pi(a_1,b,a_2) = ((a_1, b^{-1}), (b, a_2))$ and $\mu(x,y) = xy$, $\mu(a_1,b,a_2) = (a_1,a_2)$; $\mu$ is well defined since $a_1 x b^{-1} b y a_2 = a_1 xy a_2$. Together with the unit $e \in M$ and the analogous $\cX^{(3)}$ witnessing associativity, this makes $\cX$ a \emph{monoid in $\Span(\Gpd)$}.
\end{definition}

On isomorphism classes, the span gives the hyperoperation $[x] * [y] = \setof{[x b y] \st b \in U}$: the fibre of $\pi$ over $([x],[y])$ is, up to equivalence, the double coset space of the middle copy of $U$ by the parts of the isotropy of $x$ and $y$ facing the middle. The isotropy groups control the hyperproduct.

\begin{remark}
A monoid in spans of groupoids is the same as a unital $2$-Segal simplicial groupoid \cite{DK} with $X_0 = \mathrm{pt}$, $X_1 = \cX$, $X_2 = \cX^{(2)}$, $X_3 = \cX^{(3)}$; the $2$-Segal condition is associativity of the span. We do not linearize.
\end{remark}

\begin{definition}[Unit component]
$U_\cX$ is the connected component of the unit $e$. For an action groupoid as above, $U_\cX$ is the full subgroupoid on the unit group $U$. It is determined by the monoid structure and is not additional data.
\end{definition}

\subsection{The ambient category}

The two monoids of interest are $\End(W_n)$ with unit group $\Aut(W_n)$ and $\mathrm{Keller}(\AA^n)$ with unit group $\Aut(\AA^n)$. Both are $k$-points of ind-schemes: $\mathrm{Keller}(\AA^n) \subset \End(\AA^n) = \mathrm{colim}_D \setof{\text{polynomial degree} \leq D}$ is cut out by $\det JF = 1$; $\End(W_n)$ is the ind-scheme of $2n$-tuples of elements of $W_n$ (filtered by the Bernstein filtration) satisfying the Heisenberg relations; $\Aut(\AA^n)$ and $\Aut(W_n)$ are ind-groups, the latter as pairs $(\psi,\varphi)$ with $\psi\varphi = \varphi\psi = \id$. Composition is a morphism of ind-schemes. Two facts require passing from ind-schemes to a slightly weaker setting: the degree is not a morphism of ind-schemes, and the known part (\cref{known}) is an image under composition.

\begin{definition}[$\cK$]
Let $\Constr_k$ be the category of constructible subsets of $k$-points of schemes of finite type over $k$, with piecewise-regular maps (maps with constructible graph). It has finite limits, and by Chevalley's theorem (here we use that $k$ is algebraically closed) images of morphisms are constructible. Let $\cK := \Ind(\Constr_k)$. The forgetful functor $\Ind(\mathrm{Aff}_k) \to \cK$ on $k$-points preserves finite limits, so the groupoids and spans above live in $\cK$. Nonreduced structure is lost; dimension is retained.
\end{definition}

\begin{lemma}[Constructibility of $d$]
On each finite piece of $\mathrm{Keller}(\AA^n)$ the function $F \mapsto d(F)$ is constructible, and $(\NN_{>0}, \times)$, as a discrete object of $\cK$, receives a morphism $\deg$.
\end{lemma}

\begin{proof}
For an \'etale dominant map the number of points in any fibre is at most $d(F)$, with equality over a dense open. So $\setof{F : d(F) \geq m}$ is the image of the constructible set $\setof{(F, y, x_1, \dots, x_m) : x_i \text{ distinct}, F(x_i) = y}$ under projection, hence constructible, and $d(F) = \max$ fibre cardinality.
\end{proof}

On the Weyl side, constructibility of $\deg$ on each filtration piece of $\End(W_n)$ is assumed; see \cref{constrprob}.

\begin{definition}[$\fE$\label{fEdef}]
$\fE := \Mon(\Span(\Gpd(\cK))) / (\NN_{>0}, \times)$: the $(2,1)$-category whose objects are monoids in spans of ind-constructible groupoids equipped with a monoid morphism $\deg$ to the discrete monoid $(\NN_{>0}, \times)$, whose $1$-cells are strict monoid morphisms over $\NN_{>0}$, and whose $2$-cells are invertible natural transformations.
\end{definition}

Every $1$-cell of $\fE$ preserves degree; this is part of being a $1$-cell, not a property.

\begin{definition}[Keller property; reflection]
An object $\cX$ of $\fE$ has the \emph{Keller property} if $U_\cX = \deg^{-1}(1)$. A $1$-cell $f : \cX \to \cY$ \emph{reflects the unit component} if $f^{-1}(U_\cY) = U_\cX$, i.e.\ the square formed by $U_\cX \to U_\cY$ and $\cX \to \cY$ is a pullback.
\end{definition}

Always $f(U_\cX) \subseteq U_\cY$ and $U_\cY \subseteq \deg^{-1}(1)$.

\begin{lemma}[Reflection\label{refllemma}]
If the source $\cX$ has the Keller property, every $1$-cell $f : \cX \to \cY$ in $\fE$ reflects the unit component.
\end{lemma}

\begin{proof}
If $f(x) \in U_\cY$ then $\deg x = \deg f(x) = 1$, so $x \in U_\cX$.
\end{proof}

\subsection{The two vertices at fixed $n$}

\begin{definition}
\begin{enumerate}
\item $\cJ_n^\rev$: the action groupoid of $\Aut(\AA^n) \times \Aut(\AA^n)^{op}$ on $\mathrm{Keller}(\AA^n)$, with the monoid structure of \emph{reversed} composition $F \cdot F' := F' \circ F$, and $\deg = d$. It has the Keller property by \cref{keller}.
\item $\cW_n$: the action groupoid of $\Aut(W_n) \times \Aut(W_n)^{op}$ on $\End(W_n)$, with composition, and $\deg$ from \cref{degdef}. Whether it has the Keller property is \cref{qkeller}.
\end{enumerate}
\end{definition}

\begin{prop}
$L_n : \cJ_n^\rev \to \cW_n$, $F \mapsto \psi_F$ on objects and $a \mapsto \psi_a$ on automorphisms, is a $1$-cell of $\fE$. It is faithful, not full, and reflects the unit component.
\end{prop}

\begin{proof}
It is a monoid morphism by \cref{lift}(3) (reversal is why the source is $\cJ_n^\rev$), over $\NN_{>0}$ by \cref{degthm}, and ind-constructible since the adjugate of $JF$ is polynomial in the coefficients of $F$. Faithfulness is injectivity of $a \mapsto \psi_a$. It is not full since $\Aut(W_n)$ contains automorphisms (e.g.\ the exchange $q_i \leftrightarrow \partial_i$) not of the form $\psi_a$, which enlarge isotropy. Reflection is \cref{refllemma}.
\end{proof}

\subsection{Isotropy}

Each object $\psi$ of $\cW_n$ carries the isotropy group $\Iso(\psi) = \setof{(a_1, a_2) \st a_1 \psi a_2 = \psi}$, a closed sub-ind-group of $\Aut(W_n) \times \Aut(W_n)^{op}$. At the unit, $\Iso(\id) \cong \Aut(W_n)$ antidiagonally. Isotropy is an invariant of the class, strictly finer than the class; it governs the fibres of $\pi$ in \cref{spanmon}.

\begin{prop}[Isotropy of the Alp\"oge lift\label{isoprop}]
Let $c := \psi_{F_1} \in \End(W_3)$. The projection $\Iso(c) \to \Aut(W_3)^{op}$ to the right factor is injective, and its image contains the one-dimensional torus of weights $(-1,1,2)$ on $(q_1,q_2,q_3)$ and opposite weights on the $\partial_i$.
\end{prop}

\begin{proof}
Suppose $(a_1, 1) \in \Iso(c)$, i.e.\ $a_1$ fixes $\im c = k[F][D]$ pointwise. The centralizer of $k[F]$ in $W_3$ is $k[q]$: if $P \in W_3$ has order $r \geq 1$ with top symbol $\sigma(q,\xi)$, then $[P, F_i] = 0$ forces $\sum_k (\partial F_i/\partial q_k)(\partial \sigma/\partial \xi_k) = 0$ for all $i$, and invertibility of $JF$ gives $\partial\sigma/\partial\xi_k = 0$, contradicting $r \geq 1$. So $a_1$ preserves $k[q]$ and fixes $k[F]$, inducing an automorphism of $k(q)$ over $k(F)$. Since $F_1$ is not an automorphism, $k(q)/k(F)$ is not Galois (\cref{campbell}); a degree-$3$ extension that is not Galois has trivial automorphism group. So $a_1$ fixes each $q_i$, and fixes each $D_j$; since $G$ is invertible over $k[q]$, $\partial_k = \sum_j (G^{-1})_{kj} D_j$ is fixed too, so $a_1 = \id$. For the torus: the weighted homogeneity of $(a,b,c)$ means $t_\lambda \psi_{F_1} t'_\lambda{}^{-1} = \psi_{F_1}$ where $t_\lambda$ has weights $(2,1,-1)$ on the target coordinates; the left factor is itself of the form $\psi_{\text{linear}}$, so the pair lies in $\Iso(c)$ with nontrivial right component.
\end{proof}

Computing $\Iso(c)$ exactly is \cref{isoproblem}.

\subsection{Bi-invariant subsets}

The set of $\Aut(W_n)$-bi-invariant subsets of $\End(W_n)$ is the power set of $\pi_0 \cW_n$; under elementwise product it is a unital quantale, the free one on the hypermonoid $\pi_0 \cW_n$. It therefore contains no information beyond $\pi_0 \cW_n$ with its hyperoperation, and we work with $\cW_n$ itself, which also retains isotropy.

\section{Stabilization: $\FI$-indexing\label{FIsec}}

Let $\FI$ be the category of finite sets $[n] = \setof{1, \dots, n}$ and injections. For $f : [n] \mono [m]$ define
\begin{itemize}
\item on the Jacobian side, $\sigma_f(F) := F$ on the coordinates in $f([n])$ and the identity on the others, and likewise on $\Aut(\AA^n)$;
\item on the Weyl side, $\sigma_f(\psi) := \psi \otimes \id$ via $W_m = W_{f([n])} \otimes W_{[m] \sminus f([n])}$, and likewise on $\Aut(W_n)$.
\end{itemize}

\begin{prop}\label{FIprop}
\begin{enumerate}
\item $\sigma_{g \circ f} = \sigma_g \circ \sigma_f$ and $\sigma_{\id} = \id$ on the nose, so $\cJ^\rev : [n] \mapsto \cJ^\rev_n$ and $\cW : [n] \mapsto \cW_n$ are strict functors $\FI \to \fE$. Each $\sigma_f$ is faithful and not full.
\item $L_m \circ \sigma_f = \sigma_f \circ L_n$ on the nose, so $L : \cJ^\rev \Rightarrow \cW$ is a strict natural transformation, i.e.\ a $1$-cell of $\Fun(\FI, \fE)$.
\item Every $\sigma_f$ reflects the unit component. On the Jacobian side this is \cref{refllemma}. On the Weyl side: if $\psi \otimes \id$ is surjective, comparing the components at the normal-form monomial $1$ in $W_m = \bigoplus_{\mathfrak m} W_n \cdot \mathfrak m$ gives $\psi(W_n) = W_n$, and injectivity is automatic.
\item (Inner $\Sigma$-structure.) For a permutation $\pi \in \Sigma_n = \Aut_\FI([n])$, $\sigma_\pi$ is conjugation by the coordinate permutation $\pi \in \Aut(\AA^n)$, resp.\ $\in \Aut(W_n)$, and there is a $2$-cell $\theta_\pi : \sigma_\pi \Rightarrow \id$ with component $(\pi, \pi^{-1}) : x \to \pi x \pi^{-1}$ at $x$. These satisfy the cocycle condition $\theta_{\pi\rho} = \theta_\pi \circ (\sigma_\pi \theta_\rho)$, are compatible with the $\sigma_f$ via the $\FI$-group structure of $\Sigma_\bullet$, and are preserved by $L$ since $\psi_\pi$ is the same coordinate permutation on $W_n$.
\end{enumerate}
\end{prop}

By (4), on isomorphism classes all $\sigma_f$ for $f : [n] \mono [n+1]$ induce the same map $\pi_0 \cX_n \to \pi_0 \cX_{n+1}$; at the level of groupoids the separate cofaces and their coherence $2$-cells are part of the data. There is no tensor product on the vertices: $n$ is fixed at each vertex, and the only cross-$n$ structure is $\sigma$.

\section{The known part and the main diagram\label{known}}

\begin{definition}[$\cR^J$]
$\cR^J_n \subseteq \cJ^\rev_n$ is the full subgroupoid on Keller maps expressible as finite words
\begin{equation*}
a_0 \, s_1 \, a_1 \, s_2 \cdots s_m \, a_m, \qquad m \geq 0, \quad a_i \in \Aut(\AA^n), \quad s_j \in S_n,
\end{equation*}
(composition in either order; $m = 0$ gives $\Aut(\AA^n)$).
\end{definition}

\begin{prop}
$\cR^J$ is an object of $\Fun(\FI, \fE)$ and the inclusion $\iota : \cR^J \Rightarrow \cJ^\rev$ is a $1$-cell. It is componentwise a full inclusion of a sub-monoid in spans, reflects the unit component, and respects the $\theta_\pi$.
\end{prop}

\begin{proof}
Concatenation of words gives closure under the span multiplication, with $\cR^{J(2)}_n$ the part of $\cJ^{\rev(2)}_n$ over $\cR^J_n \times \cR^J_n$. Each seed family is constructible in each degree (Gallagher's parameters $p, b$ satisfy linear conditions and an open condition), so by Chevalley the words of each length form a constructible set on each finite piece. $\sigma_f$ sends words to words. The unit component of $\cR^J_n$ is $\Aut(\AA^n)$, the same as for $\cJ^\rev_n$, so $\cR^J_n$ has the Keller property and $\iota$ reflects. $\Sigma_n \subset \Aut(\AA^n)$ and $\cR^J_n$ is bi-invariant, so $\theta_\pi$ restricts.
\end{proof}

The complete situation is one diagram in $\Fun(\FI, \fE)$:
\begin{equation}\label{maindiag}
\begin{tikzcd}[column sep=5em]
\cR^J \arrow[r, Rightarrow, hook, "\iota"]
  \arrow[rr, Rightarrow, bend right=25, "L\circ\iota"']
  & \cJ^{\rev} \arrow[r, Rightarrow, "L"]
  & \cW
\end{tikzcd}
\end{equation}
All three $1$-cells reflect the unit component ($\cR^J$ has the Keller property) and respect the inner $\Sigma$-structure. The image of $L \circ \iota$ is not closed under products in $\cW$, since $\Aut(W_n)$ is larger than the image of $\Aut(\AA^n)$; we do not introduce a separate vertex for the closure.

\section{Enough}

\begin{definition}[Enough\label{enough}]
Let $I$ be a category, $X, \cY : I \to \Mon(\Span(\Gpd))$ functors, and $f : X \Rightarrow \cY$ a natural transformation. For $i \in I$ let $\cY(i)^{(m)}$ be the $m$-fold iterated multiplication object with $\pi_m : \cY(i)^{(m)} \to \cY(i)^m$ and $\mu_m : \cY(i)^{(m)} \to \cY(i)$ (defined up to equivalence by associativity; $\cY(i)^{(0)} = \mathrm{pt}$ picking out the unit), and let $S_i$ be the essential image of $f_i$. The \emph{part generated by $f$} is the subfunctor $\langle f \rangle \subseteq \cY$ with $\langle f \rangle(i)$ the full replete subgroupoid on objects $y$ such that
\begin{equation*}
y \cong \mu_m(w) \quad \text{for some } m \geq 0 \text{ and } w \in \cY(i)^{(m)} \text{ with every component of } \pi_m(w) \text{ in } S_i.
\end{equation*}
$f$ is \emph{enough} if $\langle f \rangle(i) = \cY(i)$ for every $i \in I$.
\end{definition}

Since $f_i$ preserves units and $S_i$ is replete, $S_i \supseteq U_{\cY(i)}$. $S_i$ is in general not closed under multiplication even though $f_i$ is a monoid morphism, because products in $\cY(i)$ use middle letters from all of $U_{\cY(i)}$. $\langle f \rangle$ is a subfunctor because the structure maps of $\cY$ are monoid morphisms and $f$ is natural. The case $I = \mathrm{pt}$ is the unindexed notion.

\begin{question}[Q1]\label{Q1}
Is $\iota : \cR^J \Rightarrow \cJ^\rev$ enough? Since $\cR^J_n$ is already closed under products and automorphisms, this asks whether $\iota$ is essentially surjective: is every Keller map of $\AA^n$ a word in polynomial automorphisms and coordinate-stabilized seeds?
\end{question}

\begin{question}[Q2]\label{Q2}
Is $L \circ \iota : \cR^J \Rightarrow \cW$ enough? Concretely: is every endomorphism of $W_n$ a word in automorphisms of $W_n$ and lifts of coordinate-stabilized seeds?
\end{question}

\begin{remark}[Low dimensions]
No seed lives in dimension $\leq 2$, so $\cR^J_1$ and $\cR^J_2$ consist of automorphisms. Q1 at $n = 1$ is trivially true; at $n = 2$ it is $JC_2$. Q2 at $n = 1, 2$ is exactly $DC_1, DC_2$. Anything exotic in $\cW_2$ stabilizes into the complement of $\langle L\circ\iota \rangle$ in all higher $\cW_n$.
\end{remark}

\begin{remark}[Relation between Q1 and Q2]
Q2 implies that $L : \cJ^\rev \Rightarrow \cW$ is enough. Q1 together with $L$ being enough implies Q2. Q2 does not formally imply Q1: a Keller map outside $\cR^J$ could have a lift that is generated in $\cW$ using phase-space automorphisms not of the form $\psi_a$.
\end{remark}

\section{The original conjecture and what degree decides\label{ROD}}

The first version of this document took the seed set to be $\setof{F_1}$ alone and posed, deliberately at maximal strength and expecting falsity, the \emph{Remaining Optimistic Dixmier} conjecture: that $L\circ\iota$ is enough for that seed set. Degree decides it.

\begin{prop}[Degree of reachable classes]
With seed set $\setof{F_1}$, every object of $\langle L \circ \iota \rangle(n)$ has degree $3^m$, $m \geq 0$.
\end{prop}

\begin{proof}
A word has $m$ seed letters, each of degree $3$ by \cref{degthm,degprops}(4), and automorphism letters of degree $1$; apply \cref{degprops}(2).
\end{proof}

\begin{theorem}
Assume $d(G) = 4$ for the map of \cref{deg4}. Then with seed set $\setof{F_1}$, $L\circ\iota$ is not enough: $\psi_G$ and all its coordinate stabilizations lie outside $\langle L\circ\iota \rangle(n)$ for every $n \geq 3$.
\end{theorem}

\begin{proof}
$\deg \psi_G = 4$ is not a power of $3$, and degree is constant on classes and preserved by $\sigma_f$.
\end{proof}

With the full seed set $S_3$, $\langle L\circ\iota\rangle(n)$ contains classes of every degree $d \geq 3$, and products add nothing new. Degree can then single out in $\cW_n \sminus \langle L\circ\iota\rangle(n)$ only objects of degree $1$ (excluded by the Keller property, \cref{qkeller}), degree $2$, or infinite degree (\cref{finite}). All three are expected to be empty. So, for the full seed set, degree is a grading and not a separating invariant: both Q1 and Q2 must be attacked within each degree piece. An invariant that can decide Q2 negatively must be constant on classes, preserved by products in the sense of \cref{spanmon}, and take on some object a value not attained on words; isotropy of irreducible factors and the lattice of intermediate division rings $E_\psi \subseteq M \subseteq \cD_n$ are the candidates.

\section{Open problems}

\begin{problem}[Quantum Keller\label{qkeller}]
Does $\cW_n$ have the Keller property: is every $\psi \in \End(W_n)$ with $\deg \psi = 1$ an automorphism? Equivalently, is an endomorphism of $W_n$ that extends to an automorphism of $\cD_n$ already surjective? Classically this is \cref{keller}.
\end{problem}

\begin{problem}[Finite degree\label{finite}]
Is $\deg \psi < \infty$ for every $\psi \in \End(W_n)$?
\end{problem}

\begin{problem}[Constructibility\label{constrprob}]
Is $\deg$ constructible on each Bernstein-filtration piece of $\End(W_n)$? This is needed for $\cW_n$ to be an object of $\fE$ as defined; the Jacobian side is settled.
\end{problem}

\begin{problem}[Certification of degree $4$]
Certify $d(G) = 4$ symbolically (e.g.\ by Gr\"obner elimination over $k(A,B,C)$), closing the hypothesis in \cref{ROD}.
\end{problem}

\begin{problem}[Isotropy of the seed\label{isoproblem}]
Compute $\Iso(\psi_{F_1})$ exactly. By \cref{isoprop} it embeds in $\Aut(W_3)^{op}$ and contains $\GG_m$; the question is whether the grading accounts for all of it.
\end{problem}

\begin{problem}[The product union]
Describe the fibre of $\pi$ over $([\psi_{F_1}], [\psi_{F_1}])$: is $[\psi_{F_1}] * [\psi_{F_1}] = \bigcup_a [\psi_{F_1} a \psi_{F_1}]$ a single class, finitely many, or infinitely many? Every element has degree $9$, so degree does not see this.
\end{problem}

\begin{problem}[Q1 and Q2]
\Cref{Q1,Q2}.
\end{problem}

\begin{thebibliography}{9}

\bibitem{Alpoge} L. Alp\"oge, announcement of an explicit counterexample to the Jacobian conjecture in dimension $3$, 19 July 2026.

\bibitem{Speyer} D. Speyer, \emph{The new counterexample to the Jacobian conjecture}, Secret Blogging Seminar, 20 July 2026. \url{https://sbseminar.wordpress.com/2026/07/20/the-new-counterexample-to-the-jacobian-conjecture/}

\bibitem{Tao} T. Tao, \emph{A digestion of the Jacobian conjecture counterexample}, 21 July 2026. \url{https://terrytao.wordpress.com/2026/07/21/a-digestion-of-the-jacobian-conjecture-counterexample/}

\bibitem{fam} A. Gallagher, \emph{Jacobian explained}, \url{https://jacobianfun.org/jacobian-explained}; see also D. Speyer's note \url{https://sbseminar.wordpress.com/wp-content/uploads/2026/07/jacobiantangentsweep.pdf}. Unrefereed; cited for the construction, with the properties we rely on marked where used.

\bibitem{BCW} H. Bass, E. H. Connell, D. Wright, \emph{The Jacobian conjecture: reduction of degree and formal expansion of the inverse}, Bull. Amer. Math. Soc. 7 (1982), 287--330.

\bibitem{Campbell} L. A. Campbell, \emph{A condition for a polynomial map to be invertible}, Math. Ann. 205 (1973), 243--248.

\bibitem{CvdD} E. H. Connell, L. van den Dries, \emph{Injective polynomial maps and the Jacobian conjecture}, J. Pure Appl. Algebra 28 (1983), 235--239.

\bibitem{BKK} A. Belov-Kanel, M. Kontsevich, \emph{The Jacobian conjecture is stably equivalent to the Dixmier conjecture}, Mosc. Math. J. 7 (2007), 209--218.

\bibitem{DK} T. Dyckerhoff, M. Kapranov, \emph{Higher Segal spaces}, Lecture Notes in Math. 2244, Springer, 2019.

\bibitem{vdE} A. van den Essen, \emph{Polynomial Automorphisms and the Jacobian Conjecture}, Progress in Mathematics 190, Birkh\"auser, 2000.

\end{thebibliography}

\end{document}
